\section{Цели экспериментальной валидации}

Данный раздел описывает план экспериментальной проверки архитектуры DREAM-Net. Все представленные в предыдущих разделах механизмы являются гипотезами, требующими эмпирической верификации.

\textbf{Основные цели:}
\begin{enumerate}
    \item Подтвердить работоспособность архитектуры на задачах обработки последовательностей.
    \item Оценить преимущество онлайн-адаптации перед статическими архитектурами.
    \item Проверить устойчивость системы к нестационарным данным и шуму.
    \item Количественно измерить влияние отдельных компонентов через абляционные тесты.
\end{enumerate}

\textbf{Важное замечание:} Результаты, представленные в черновиках спецификации (например, loss 0.0024 на LJSpeech), являются предварительными и требуют воспроизведения в контролируемых условиях с указанием дисперсии и статистической значимости.

\section{Задачи и датасеты}

\subsection{Базовая задача ASR (Automatic Speech Recognition)}

\textbf{Датасет:} LJSpeech (английский, один диктор)

\begin{table}[h]
\centering
\caption{Параметры датасета LJSpeech}
\label{tab:ljspeech}
\begin{tabular}{lp{8cm}}
\toprule
\textbf{Параметр} & \textbf{Значение} \\
\midrule
Объём & $\sim$24 часа речи \\
Частота дискретизации & 22.05 кГц \\
Разметка & Текст на уровне предложений \\
\bottomrule
\end{tabular}
\end{table}

\textbf{Цель:} Базовая валидация сходимости модели на стандартном датасете.

\textbf{Метрики:}
\begin{itemize}
    \item Character Error Rate (CER)
    \item Word Error Rate (WER)
    \item Функция потерь (Loss) по эпохам
\end{itemize}

\subsection{Задача адаптации к диктору}

\textbf{Датасет:} Mozilla Common Voice (множество дикторов)

\begin{table}[h]
\centering
\caption{Параметры датасета Common Voice}
\label{tab:commonvoice}
\begin{tabular}{lp{8cm}}
\toprule
\textbf{Параметр} & \textbf{Значение} \\
\midrule
Количество дикторов & $\geq 100$ \\
Языки & Английский, Русский (опционально) \\
Разметка & Текст + ID диктора \\
\bottomrule
\end{tabular}
\end{table}

\textbf{Протокол теста «Смена диктора»:}
\begin{enumerate}
    \item Модель обучается на дикторах A, B, C.
    \item Во время инференса происходит переключение на диктора D (не виден при обучении).
    \item Измеряется количество шагов до восстановления базовой точности.
\end{enumerate}

\textbf{Гипотеза:} DREAM-Net восстановит точность за $< 50$ шагов благодаря быстрым весам $U$, тогда как статические модели потребуют fine-tuning.

\subsection{Задача робастности к шуму}

\textbf{Датасет:} LJSpeech + искусственный шум

\begin{table}[h]
\centering
\caption{Типы и интенсивность шума}
\label{tab:noise}
\begin{tabular}{lp{6cm}}
\toprule
\textbf{Тип шума} & \textbf{Интенсивность (SNR)} \\
\midrule
Белый шум & 0 dB, 10 dB, 20 dB \\
Розовый шум & 0 dB, 10 dB, 20 dB \\
Реверберация & T60 = 0.3s, 0.6s, 0.9s \\
Фоновая речь & 1 говорящий, 3 говорящих \\
\bottomrule
\end{tabular}
\end{table}

\textbf{Цель:} Проверить эффективность Surprise Gate в фильтрации аномалий.

\textbf{Метрика:} CER в зависимости от SNR (кривая деградации).

\section{Базовые архитектуры для сравнения (Baselines)}

Для корректной оценки необходимо сравнение с современными архитектурами аналогичной ёмкости.

\begin{table}[h]
\centering
\caption{Базовые архитектуры для сравнения}
\label{tab:baselines}
\begin{tabular}{lcp{6cm}}
\toprule
\textbf{Архитектура} & \textbf{Параметры} & \textbf{Обоснование выбора} \\
\midrule
\textbf{LSTM} & $\sim$100k & Классическая RNN-базлайн \\
\textbf{GRU} & $\sim$100k & Упрощённая рекуррентная архитектура \\
\textbf{Transformer (Small)} & $\sim$500k & Attention-based, квадратичная сложность \\
\textbf{Mamba / SSM} & $\sim$100k & Линейная сложность, состояние \\
\textbf{CfC (Liquid Networks)} & $\sim$100k & Непрерывная динамика, биовдохновлённая \\
\textbf{DREAM-Net} & $\sim$82k & Предлагаемая архитектура \\
\bottomrule
\end{tabular}
\end{table}

\textbf{Важно:} Все модели должны быть обучены на одинаковых данных, с одинаковой предобработкой и аугментацией для честного сравнения.

\section{Ключевые метрики оценки}

\subsection{Метрики точности}

\begin{table}[h]
\centering
\caption{Метрики точности}
\label{tab:accuracy_metrics}
\begin{tabular}{lp{6cm}p{4cm}}
\toprule
\textbf{Метрика} & \textbf{Формула} & \textbf{Назначение} \\
\midrule
\textbf{CER} & $\frac{S + D + I}{N} \times 100\%$ & Точность на уровне символов \\
\textbf{WER} & $\frac{S + D + I}{N} \times 100\%$ & Точность на уровне слов \\
\textbf{Loss} & $\mathcal{L}_{total}$ & Сходимость обучения \\
\bottomrule
\end{tabular}
\end{table}
где $S$ — замены, $D$ — удаления, $I$ — вставки, $N$ — общее количество единиц.

\subsection{Метрики адаптивности}

\begin{table}[h]
\centering
\caption{Метрики адаптивности}
\label{tab:adaptivity_metrics}
\begin{tabular}{lp{6cm}l}
\toprule
\textbf{Метрика} & \textbf{Описание} & \textbf{Единица} \\
\midrule
\textbf{Adaptation Time} & Шагов до восстановления CER после смены диктора & шаги \\
\textbf{Surprise Correlation} & Корреляция $S_t$ с реальными аномалиями & коэффициент Пирсона \\
\textbf{Plasticity Ratio} & Доля шагов с активным обновлением $U$ & процент \\
\bottomrule
\end{tabular}
\end{table}

\subsection{Метрики эффективности}

\begin{table}[h]
\centering
\caption{Метрики эффективности}
\label{tab:efficiency_metrics}
\begin{tabular}{lp{6cm}l}
\toprule
\textbf{Метрика} & \textbf{Описание} & \textbf{Единица} \\
\midrule
\textbf{Parameters} & Общее количество обучаемых параметров & count \\
\textbf{FLOPs/step} & Операций на один шаг инференса & FLOPs \\
\textbf{Memory/step} & Памяти на один шаг (включая контекст) & MB \\
\textbf{Real-time Factor} & Время обработки / длительность аудио & коэффициент \\
\bottomrule
\end{tabular}
\end{table}

\subsection{Метрики стабильности}

\begin{table}[h]
\centering
\caption{Метрики стабильности}
\label{tab:stability_metrics}
\begin{tabular}{lp{6cm}l}
\toprule
\textbf{Метрика} & \textbf{Описание} & \textbf{Единица} \\
\midrule
\textbf{Forgetting Rate} & Деградация CER на старых дикторах после сна & \% \\
\textbf{Shock Frequency} & Среднее количество входов в Shock per минуту & count/min \\
\textbf{Weight Drift} & $\|U_t - U_{target}\|_F$ за сессию & норма \\
\bottomrule
\end{tabular}
\end{table}

\section{Абляционные тесты}

Для оценки вклада отдельных компонентов архитектуры предлагается следующая серия абляционных тестов:

\begin{table}[h]
\centering
\caption{Абляционные тесты}
\label{tab:ablation}
\begin{tabular}{lp{5cm}p{5cm}}
\toprule
\textbf{Тест} & \textbf{Модификация} & \textbf{Гипотеза} \\
\midrule
\textbf{A-1} & Без Surprise Gate ($S_t = 1$ всегда) & Увеличится переобучение на шуме \\
\textbf{A-2} & Без фазы сна (нет консолидации) & Увеличится forgetting rate \\
\textbf{A-3} & Без быстрых весов ($U_t = 0$) & Снизится адаптивность к диктору \\
\textbf{A-4} & Без контекста ($c_t = 0$) & Ухудшится работа на длинных последовательностях \\
\textbf{A-5} & Без гистерезиса (мгновенные переходы) & Увеличится дребезг режимов \\
\textbf{A-6} & Без Ramp-up (мгновенное восстановление) & Возможны вторичные шоки \\
\textbf{A-7} & Без Double Buffering (синхронный сон) & Блокировки инференса во время сна \\
\textbf{A-8} & Без Clipping весов & Риск численной нестабильности \\
\bottomrule
\end{tabular}
\end{table}

\textbf{Протокол:} Каждый тест запускается на полном наборе задач (6.2.1–6.2.3). Результаты сравниваются с полной версией DREAM-Net.

\section{Протокол валидации}

\subsection{Разделение данных}

\begin{table}[h]
\centering
\caption{Разделение данных}
\label{tab:data_split}
\begin{tabular}{lcp{6cm}}
\toprule
\textbf{Набор} & \textbf{Доля} & \textbf{Назначение} \\
\midrule
\textbf{Train} & 70\% & Обучение медленных весов $W$ \\
\textbf{Validation} & 15\% & Подбор гиперпараметров \\
\textbf{Test} & 15\% & Финальная оценка (не используется для настройки) \\
\bottomrule
\end{tabular}
\end{table}

\subsection{Кросс-валидация}

Для оценки статистической значимости результатов рекомендуется k-fold кросс-валидация:
\begin{itemize}
    \item $k = 5$ (пять фолдов)
    \item Для каждого фолда: отдельный запуск обучения с разными seed
    \item Отчёт: среднее значение $\pm$ стандартное отклонение
\end{itemize}

\subsection{Критерии успеха}

Архитектура считается успешно валидированной при выполнении следующих условий:

\begin{table}[h]
\centering
\caption{Критерии успеха валидации}
\label{tab:success_criteria}
\begin{tabular}{lp{5cm}l}
\toprule
\textbf{Критерий} & \textbf{Порог} & \textbf{Приоритет} \\
\midrule
CER на LJSpeech & $\leq 10\%$ & Высокий \\
Adaptation Time & $< 100$ шагов & Высокий \\
Forgetting Rate & $< 5\%$ после 10 циклов сна & Средний \\
Real-time Factor & $< 1.0$ (быстрее реального времени) & Средний \\
Численная стабильность & 0 NaN/Inf за 100 часов & Критический \\
\bottomrule
\end{tabular}
\end{table}

\textbf{Примечание:} Пороги являются предварительными и могут быть скорректированы по результатам первых экспериментов.

\section{План отчётности}

Для каждого эксперимента должен быть подготовлен отчёт со следующей структурой:

\begin{enumerate}
    \item \textbf{Конфигурация:} Версия кода, гиперпараметры, seed, оборудование.
    \item \textbf{Данные:} Версия датасета, предобработка, аугментация.
    \item \textbf{Результаты:} Таблицы с метриками (среднее $\pm$ std), графики обучения.
    \item \textbf{Анализ:} Интерпретация результатов, выявленные аномалии.
    \item \textbf{Выводы:} Подтверждена ли гипотеза, рекомендации для следующей итерации.
\end{enumerate}

\textbf{Требование воспроизводимости:} Весь код, конфигурации и предобученные веса должны быть зафиксированы в системе контроля версий с возможностью воспроизведения результатов.

\section{Ограничения экспериментального плана}

Следующие аспекты требуют дополнительного уточнения:

\begin{enumerate}
    \item \textbf{Вычислительные ресурсы:} Полный набор экспериментов может потребовать значительных GPU-часов. Рекомендуется приоритизация задач.
    
    \item \textbf{Доступ к данным:} Некоторые датасеты (например, Common Voice) требуют регистрации и могут иметь ограничения на использование.
    
    \item \textbf{Время валидации:} Тесты на долгосрочное забывание (10+ циклов сна) требуют длительных прогонов.
    
    \item \textbf{Статистическая мощность:} Для надёжной оценки значимости различий может потребоваться больше чем 5 фолдов.
\end{enumerate}

Эти ограничения следует учитывать при планировании экспериментальной фазы.

\section{Сводная таблица экспериментов}

\begin{table}[h]
\centering
\caption{План экспериментов}
\label{tab:experiments_summary}
\begin{tabular}{lp{3cm}p{3cm}lp{2cm}}
\toprule
\textbf{Эксперимент} & \textbf{Датасет} & \textbf{Метрики} & \textbf{Базлайны} & \textbf{Приоритет} \\
\midrule
E-1: Базовый ASR & LJSpeech & CER, WER, Loss & LSTM, CfC, Mamba & Высокий \\
E-2: Адаптация & Common Voice & Adaptation Time & Transformer, SSM & Высокий \\
E-3: Шум & LJSpeech + Noise & CER vs SNR & Все базлайны & Средний \\
E-4: Абляция & LJSpeech & Все метрики & DREAM-Net full & Средний \\
E-5: Долгосрочность & Custom & Forgetting Rate & Модели без сна & Низкий \\
\bottomrule
\end{tabular}
\end{table}
